% lesson-peel.tex -- shared lesson content (no styling here).
% Uses only the semantic macros documented in ../MACROS.md.

\title{Building a body paragraph with PEEL}
\subtitle{Argument structure for Writing Task 2}
\author{Lin Chen}
\examtag{IELTS Academic \textperiodcentered{} Writing Task 2}
\lessonlabel{Lesson 7}
\duration{60 min}
\target{Band 6.0 to 6.5}
\learner{Jordan}
\date{}

\begin{document}

\titleslide

% ---------------------------------------------------------------- 2
\begin{frame}[stage=Goals]{Today's objectives}
  \begin{objectives}
    \item Name the four PEEL moves and what each one does.
    \item Turn a vague opinion into one precise, arguable Point sentence.
    \item Write a 120--150-word body paragraph in eight minutes.
  \end{objectives}
  \bandfocus{Task Response \textperiodcentered{} Coherence and Cohesion}
\end{frame}

% ---------------------------------------------------------------- 3
\begin{frame}[stage=Diagnose]{From last week's essay}
  \begin{compare}
    \original{Nowadays, many people \err{thinks} that \err{study other subject} is
      \err{waste of time}, but I think it \err{have many advantage}.}
    \revised{Requiring students to take courses outside their major \fix{has} one clear
      advantage: it develops the \fix{flexible thinking} that employers value.}
    \why{\emph{many people think}: plural subject, plural verb}
    \why{\emph{studying other subjects}: a gerund as subject}
    \why{\emph{a waste of time}: the noun needs an article}
    \why{\emph{it has}, and name the advantage you mean}
  \end{compare}
\end{frame}

% ---------------------------------------------------------------- 4
\begin{frame}[stage=Structure]{Four moves, one argument}
  \peel{State one clear claim that answers the question.}
       {Give the reason the claim is true.}
       {Make it concrete with a specific case.}
       {Return to the question and your position.}
  \aside{Most Band 6 paragraphs jump from Point straight to Example. Explain is the move that lifts the band.}
\end{frame}

% ---------------------------------------------------------------- 5
\begin{frame}[stage=Speak first]{Say it before you write it}
  \begin{cuecard}{Describe a subject you studied that was not part of your main interest.}
    \cue{what the subject was}
    \cue{when and where you studied it}
    \cue{what you learned from it}
    \cueexplain{and explain whether you think it was worth studying.}
    \cuetiming{1 min to prepare \textperiodcentered{} 1--2 min to speak}
  \end{cuecard}
  \aside{Jordan speaks first. Then they write the same idea as one PEEL paragraph.}
\end{frame}

% ---------------------------------------------------------------- 6
\begin{frame}[stage=Practise]{Write one body paragraph}
  \begin{timed}{8:00}
    \prompt{Some people believe that university students should be required to take
      courses outside their main subject. To what extent do you agree or disagree?}
    \begin{phases}
      \phase{0--1 min}{Plan: choose one Point and one example.}
      \phase{1--6 min}{Write 120--150 words in PEEL order.}
      \phase{6--8 min}{Check verb agreement, articles and the Link.}
    \end{phases}
  \end{timed}
\end{frame}

% ---------------------------------------------------------------- 7 (overlays)
\begin{frame}[stage=Model]{A model paragraph, move by move}
  \begin{reveal}
    \step[One claim, and no examples yet.]{Point}{The strongest argument for compulsory
      courses outside a student's major is that they build transferable thinking skills.}
    \step[The reasoning: why the claim holds.]{Explain}{Unfamiliar subjects force students
      to question their assumptions instead of applying familiar methods automatically.}
    \step[Specific and believable, not \emph{many people}.]{Example}{For instance, a computer
      science student who takes an ethics module learns to ask not only whether software
      works, but whom it might harm.}
    \step[Returns to the question and the writer's position.]{Link}{Broad requirements
      therefore produce more adaptable graduates, which is why I largely agree with this policy.}
  \end{reveal}
\end{frame}

% ---------------------------------------------------------------- 8
\keyframe[Rule of thumb]{One paragraph, one point.\\Every other sentence serves it.}

% ---------------------------------------------------------------- 9
\begin{frame}[stage=Wrap-up]{Summary and exit check}
  \begin{columns}[T,onlytextwidth]
    \column{0.45\linewidth}
      \begin{summary}
        \item A Point makes one specific claim.
        \item Explain gives the reason; Example makes it concrete.
        \item The Link answers the question, not the topic.
      \end{summary}
      \begin{wordlist}
        \vocab{transferable}{trænsˈfɜːrəbl}{useful in many situations}
        \vocab{compulsory}{kəmˈpʌlsəri}{required by rule}
      \end{wordlist}
    \column{0.49\linewidth}
      \begin{exitcheck}
        \item Which move is missing?
          \sample{Online study saves time. For example, I save two hours a day.}
        \item Correct this sentence: \sample{It have many advantage.}
        \item Write a Point sentence for: \sample{Should museums be free?}
      \end{exitcheck}
  \end{columns}
\end{frame}

\end{document}
